% Part of main.tex; compile main.tex.
% ---------------------------------------------------------------------------
\section{Conclusions and Recommendations}
\label{sec:conclusion}

Sleep-based compute emulation answers a narrow question: \emph{if the
application left every node resource idle during its computation phase, how
much of the I/O could be hidden?} This is a useful upper bound, but it is
not the question practitioners need answered. In this paper, we surveyed 17
evaluations and benchmarks of asynchronous I/O and found that \texttt{sleep},
an unspecified compute phase, or no compute phase at all is the norm. We then
measured five asynchronous I/O implementations with compute surrogates that
consume resources. With these surrogates, the application ran longer than a
\texttt{sleep}-based evaluation predicts in \ExcLongerFive of \ExcN
configurations, by up to \ExcMaxPct\%. Depending on the implementation, the
cost appeared as a slower computation or as I/O that did not progress during
the computation. Based on these results and on the surveyed literature, we
provide the following recommendations.

\paragraph{For authors evaluating asynchronous I/O.}
\begin{enumerate}
  \item \textbf{Treat sleep as an upper bound, and label it as such.} Results
        obtained with \texttt{sleep()} should be presented as the best-case
        overlap, not as the expected benefit for an application.
  \item \textbf{Use compute surrogates that consume resources.} At a minimum,
        include a CPU-bound and a memory-bandwidth-bound surrogate calibrated
        to the same duration, as overlap benchmarks such as IMB-IO already
        do~\cite{intel_imb_io}. In addition, include a communication-bound
        surrogate, which was the most sensitive surrogate in our
        measurements.
  \item \textbf{Report the end-to-end time and attribute it.} The figure of
        merit is the run time that the application observes. Report it, and
        break the difference from the ideal into compute-phase slowdown and
        exposed I/O. An I/O time that drops to near zero while the compute
        phase grows is not an improvement, and neither is an undisturbed
        compute phase followed by a long wait for I/O that never progressed.
        A sleep phase can reveal neither case.
  \item \textbf{Report the core allocation explicitly.} State whether cores
        were reserved for background threads and how many ranks and threads
        ran per node. In addition, evaluate both fully subscribed and
        reserved configurations, because the placement changed our results
        considerably.
  \item \textbf{Validate on at least one real application.} Kernel results
        should be checked against a full application. The gap between
        sleep-based I/O-time speedups and end-to-end speedups in real codes,
        such as Flash-X~\cite{jain2022flashx}, shows why.
  \item \textbf{Do not tune the compute phase to hide the I/O.} Choosing the
        emulated compute length so that the overlap is complete guarantees
        the outcome. Instead, derive $t_{comp}$ from measured application
        profiles, and sweep it to show where the benefit begins and ends.
  \item \textbf{Repeat the measurements, and check the synchronous
        baseline.} I/O time can vary considerably between runs, as it did on
        our node, and a single synchronous run can inflate or deflate a
        speedup. In addition, \texttt{sleep} itself slowed down the
        synchronous I/O in our measurements, by \RecheckMinPct\% to
        \RecheckMaxPct\%, which inflates \texttt{sleep}-based speedups
        further. Comparing asynchronous runs directly, as we do, reduces
        this sensitivity.
  \item \textbf{Benchmark suites should offer more than sleep.} Tools such as
        h5bench~\cite{li2021h5bench} could provide selectable compute
        surrogates alongside the existing emulated-compute parameter.
\end{enumerate}

\paragraph{For readers interpreting existing results.}
\begin{enumerate}
  \item \textbf{Read sleep-based speedups as upper bounds.} They show the
        overhead of the asynchronous mechanism without contention and the
        maximum attainable hiding, which is valuable, but they do not
        predict the speedup of an application.
  \item \textbf{Check what is being measured.} Many studies report the
        observed I/O time with the emulated computation excluded. Such
        numbers cannot reveal a slower compute phase or I/O that stalls
        behind the computation, and they should not be equated with an
        end-to-end speedup.
  \item \textbf{Consider the resource profile of the workload.} Sleep is a
        closer approximation when the host CPU is actually idle during the
        computation, as for GPU-offloaded codes~\cite{tang2022async}. In our
        measurements with pinned ranks, \texttt{sleep} predicted the
        asynchronous run time with a GPU kernel within a few percent for
        POSIX, HDF5, and MPI-IO, but not for ADIOS2. Even in this case, the
        speedup should be computed against a synchronous run with the real
        computation, because \texttt{sleep} slows down the synchronous
        baseline. Sleep is a poor approximation for CPU-bound,
        memory-bandwidth-bound, or communication-heavy codes running on fully
        subscribed nodes.
  \item \textbf{Look for stated caveats.} Statements such as ``in the absence
        of node-level oversubscription'' or recommendations to reserve a core
        mark the boundary of the validity of a result.
\end{enumerate}

Asynchronous I/O remains a promising technique, and the studies surveyed in
this paper show that its mechanisms can be made efficient. Our argument is
about the evidence, not about the technique. Ultimately, to know how much
asynchronous I/O helps a real application, its evaluation must include a
computation phase that actually computes.
